% Step 1.1: synthetic benchmark, geometry and ground-truth flow / FFR.
% Numbers from results/step1.1_summary.json.
% Bib keys to add to the .bib: zeng2023imagecas, thakur2026punch, young1973flow.

\paragraph{Vessel geometry.}
No public dataset pairs dynamic angiography with ground-truth flow or FFR,
so we build a synthetic benchmark on real coronary anatomy. From the
expert coronary labels of ImageCAS~\cite{zeng2023imagecas} we separate the
left coronary tree, skeletonize it on an isotropic grid at the native
in-plane spacing, and take the radius from the Euclidean distance
transform. Because ImageCAS does not label individual arteries, the LAD is
selected as the most anterior long root-to-leaf path from the ostium, and
every selection passes an automatic and a visual quality check. The
centerline and radius are smoothed and resampled every 0.5\,mm. The
resulting 40 LADs measure 81 to 165\,mm from the ostium.

\paragraph{Stenoses, flow and FFR.}
Each LAD receives five independent focal stenoses with a smooth cosine
profile, a diameter reduction drawn uniformly from 30 to 70\,\% and a
length from 5 to 15\,mm, placed in the proximal 60\,\% of the vessel,
giving 200 cases. Following PUNCH~\cite{thakur2026punch}, the mean resting
velocity in the normal segment is drawn from 10 to 20\,cm/s and the
hyperemic flow is the resting flow scaled by a coronary flow reserve drawn
from 1.5 to 4.0, capped at 0.4\,m/s. Ground-truth FFR is computed at
hyperemic flow with the Young--Tsai pressure-drop model~\cite{young1973flow}
($\mu = 3.5$\,mPa\,s, $\rho = 1060$\,kg/m$^3$, $K_t = 1.52$,
$P_a = 90$\,mmHg). The benchmark spans FFR values from 0.30 to 0.99, with
36\,\% of cases below the 0.80 treatment threshold. At rest, the contrast
transit time through the LAD ranges from 0.28 to 2.0\,s, so even with 30
views over the 4.4\,s rotational sweep the bolus front is observed in at
least two consecutive views.
